\documentclass[11pt]{article}
\usepackage{fontspec,amsmath,amssymb,amsthm,hyperref}
\usepackage[margin=1in]{geometry}
\setlength{\emergencystretch}{3em}
\title{Uniform weak consistency for compact smooth tests}
\author{LangevinLean mathematical exposition}
\date{}
\begin{document}\maketitle
\noindent Source module: \texttt{AverageLMC/LaplacianWeakStep.lean}. Compile this proof document with XeLaTeX or LuaLaTeX. The accompanying HTML pairs the mathematical argument with the exact Lean statements and records the scope of the formalization.


This smooth-test development supplies the Gaussian calculus and drift-noise separation used by the later C\ensuremath{{}^{2}} extension. Smoothness is imposed only on auxiliary compact tests; the final potential keeps the paper's C\ensuremath{{}^{3}} hypothesis.



Drift and noise are separated by the deterministic estimate


\[\left|\phi(x+w+d)-\phi(x+w)-D\phi(x)[d]\right|
\le\frac L2\|d\|^2+L\|w\|\,\|d\|.\]


For every compact smooth test there is a constant \(K\ge0\), uniform in \(x,v\) and \(0\le h\le1\), such that


\[\left|\mathbb E\phi\bigl(x+hv+\sqrt{2h}Z\bigr)-\phi(x)
-h\bigl(D\phi(x)[v]+\Delta\phi(x)\bigr)\right|
\le Kh\sqrt h\,(1+\|v\|^2).\]
\section{Bound the test gradient globally}


Compact support bounds the second derivative everywhere. The derivative norm bound gives a Lipschitz derivative, and the Riesz isometry transfers this constant to the gradient.



\noindent Lean declarations: \texttt{compactSmooth\_gradient\_lipschitz}.

\section{Control a deterministic drift increment}


Apply the smooth upper bound to φ and −φ to obtain an absolute first-order remainder bounded by L‖d‖\ensuremath{{}^{2}}/2. Apply it at x+w and compare the two derivatives by gradient Lipschitz continuity, obtaining the additional L‖w‖‖d‖ term. This separates a potentially large random drift from the small Gaussian increment.



\noindent Lean declarations: \texttt{firstOrder\_remainder\_abs\_le}, \texttt{shifted\_firstOrder\_remainder\_abs\_le}.

\section{Average the true Gaussian update}


Use the compact smooth Taylor remainder for the noise alone and the shifted first-order estimate for the drift. Gaussian first and third moments are finite. For 0\ensuremath{\le}h\ensuremath{\le}1, absorb the deterministic h\ensuremath{{}^{2}} term into h√h and obtain one uniform constant K times h√h(1+‖v‖\ensuremath{{}^{2}}), for arbitrary starting point and drift.



\noindent Lean declarations: \texttt{compactSmooth\_gaussian\_weak\_step}.

\section{Identify regularity of the diffusion term}


Express the actual Laplacian as a finite sum of second derivative evaluations on an orthonormal basis. Continuity of the second derivative proves continuity of this sum.



\noindent Lean declarations: \texttt{compactSmooth\_continuous\_laplacian}.

\section{Insert an actual bounded past weight}


Integrate the pointwise Gaussian weak residual against a measurable weight of absolute value at most A. The drift needs only its true second moment, producing A K h√h(1+E‖v‖\ensuremath{{}^{2}}). Product measurability and this integrable bound justify the averaging. The separate C\ensuremath{{}^{2}} module refines this API to the regularity needed by the gradient cutoffs.



\noindent Lean declarations: \texttt{compactSmooth\_gaussian\_weighted\_weak\_step}.

\end{document}
